\section{Video generation 与研究议程：把预测写成可验证问题}

视频把前面所有瓶颈叠加：token/latent 更多、temporal consistency 更难、high-resolution 更贵、data coverage 与 caption quality 更差。讲者没有给出 Sora 的内部结论，而是从公开现象拆出一套 2024 年的工程假设，并据此提出研究方向。

\subsection{Sora-like recipe：3D latents、scale、context 与 data}

课堂以 CogVideo 为早期 autoregressive text-to-video 参照，再把 Sora-like improvement 分成 deflickering、image quality、resolution、long-context infrastructure 与 recaption/coverage。若视频 latent 大小为 $T\times H\times W$，naive token 数为
\[
n_v=T\frac{H}{p_h}\frac{W}{p_w},
\]
full attention 的二次成本会迅速爆炸，因此 3D latent compression、factorized attention、parallelism 与 data pipeline 都是必要条件。

\lecturefigure{slide-29-video-generation.jpg}{Video generation：从 CogVideo 到 Sora-like scaling recipe}{Stanford Online 官方录像 00:58:29}

\begin{knowledgebox}{五个工程因素分别解决什么}
\small\sloppy
\begin{tabularx}{\textwidth}{L{0.19\textwidth} L{0.32\textwidth} X}
\toprule
因素 & 主要问题 & 不能自动保证的能力 \\
\midrule
3D latent encoder & 压缩时空并减少 flicker & 真实物理与长期因果 \\
Model/data scale & 提升视觉质量和覆盖 & 无偏、安全、可控 \\
High resolution & 保存文字与细节 & 正确动作与几何 \\
Long-context infra & 承载更多 frames/tokens & 有效利用所有 context \\
Recaption 与 coverage & 改善 text-video supervision & 从 raw video 学到物理规律 \\
\bottomrule
\end{tabularx}
\end{knowledgebox}

\teachervoice{00:57:54--01:04:00，讲者把这些因素称为对 Sora-like progress 的总结，不是 controlled ablation。讲义因此把它们写成可检验 hypotheses，而不是已知内部配方。}

\subsection{2024 年预测：识别、视频理解与 embodied AI}

前一节拆解的是模型和数据怎样扩大 video generation；这里转向更不确定的 research forecast。预测 slide 认为高层视觉识别会更便宜，video understanding 会成为新热点，embodied AI 会产生亮眼 demo 但短期难进入日常生活。正确记录方式是保留 2024 年 5 月的时间、资源假设与“可能”语气，而不是事后只挑选命中的句子。

\lecturefigure{slide-30-research-predictions.jpg}{2024 年课堂预测：视觉识别、video understanding 与 embodied AI}{Stanford Online 官方录像 01:04:00}

\begin{warningbox}{Prediction 不是 benchmark result}
“一两年”“80\% solved”“无法很快影响现实生活”都不是可直接测量的标准。要验收预测，必须重新定义 task distribution、cost、accuracy、long-tail、deployment environment 与时间点。
\end{warningbox}

\teachervoice{01:00:00--01:04:00，讲者同时表达乐观与限制：视觉识别会快速便宜化，但 video model 仍有 hallucination/counting 问题；embodied AI 会有惊艳 demo，却可能昂贵且难以日常部署。}

\subsection{研究建议：按资源与目标选择问题}

最后一页不是“选题排行榜”，而是把 career objective、compute access 和 research risk 联系起来。video datasets/benchmarks 适合资源较少但能组织 evaluation 的团队；speech 被认为需求高但投入不足；architecture/optimizer 可能影响大但要求硬件与系统能力；compute-to-data 则是 web corpus 逐渐饱和后的核心方向。

\lecturefigure{slide-31-research-advice.jpg}{按目标与资源选择研究问题：video、speech、systems、architecture 与 data}{Stanford Online 官方录像 01:07:34}

\begin{knowledgebox}{把建议改写成决策矩阵}
\small\sloppy
\begin{tabularx}{\textwidth}{L{0.17\textwidth} L{0.20\textwidth} L{0.20\textwidth} X}
\toprule
方向 & 主要资产 & 主要风险 & 可验证交付 \\
\midrule
视频数据/评测 & collection、annotation、evaluation & 泄漏、shortcut、覆盖不足 & dataset card、held-out test、error taxonomy \\
Speech/audio & 用户场景与多语言数据 & 隐私、噪声、实时延迟 & ASR、TTS 与对话端到端指标 \\
系统协作 & GPU 拓扑与 kernel 知识 & 复现与硬件依赖 & 吞吐、延迟、利用率与成本 \\
架构/优化器 & 大规模 controlled runs & 计算昂贵、结论不稳 & matched-compute scaling/ablation \\
Compute-to-data & verifier、execution、search & feedback bias、synthetic collapse & provenance、verification rate、novelty \\
\bottomrule
\end{tabularx}
\end{knowledgebox}

\begin{lstlisting}[caption={多模态研究项目的发布前检查}]
freeze_task_definition_and_data_split()
record_model_version_resolution_and_prompt()
report_compute_latency_memory_and_cost()
separate_demo_success_from_repeated_trial_rate()
audit_shortcuts_leakage_and_long_tail_failures()
publish_provenance_metrics_and_negative_results()
\end{lstlisting}

\teachervoice{01:05:00--01:08:10，讲者称 speech AI 被低估，并建议做算法的人尽早与 systems researchers 合作，因为最好的算法必须映射到当代硬件。}

\subsection{本章小结}

Video generation 不是把 image model 多跑几帧，而是时空压缩、并行训练、long context、数据覆盖和评价共同组成的系统问题。研究建议也应按资源与可验证交付理解，而不是把课堂预测当成确定路线图。

